%! Author = lili
%! Date = 7/25/26


\newglossaryentry{def_komplement}{
    name={Wie ist das Komplement $A^c$ einer Menge $A\subseteq\Omega$ definiert?},
description={$A^c := \Omega \setminus A$. Achtung: Diese Notation setzt voraus, dass klar ist, auf welche Referenz-Obermenge $\Omega$ man sich bezieht.},
type=dinge
}

\newglossaryentry{def_vereinigung_schnitt_abzaehlbar}{
name={Wie sind Vereinigung und Schnitt abzählbar vieler Mengen $A_i$, $i\in I$, definiert?},
description={$\bigcup_{i\in I} A_i = \{\omega\in\Omega : \exists j\in I \text{ mit } \omega\in A_j\}$ (mind. ein $A_i$ enthält $\omega$). $\bigcap_{i\in I} A_i = \{\omega\in\Omega : \omega\in A_i \ \forall i\in I\}$ (alle $A_i$ enthalten $\omega$).},
type=dinge
}

\newglossaryentry{lem_rechenregeln_mengen}{
name={Welche drei Rechenregeln für Mengen/Ereignisse kennst du (Kommutativ-, Assoziativ-, de Morgan)?},
description={Kommutativität: $A\cup B = B\cup A$, $A\cap B = B\cap A$. Assoziativität: $(A\cup B)\cup C = A\cup(B\cup C)$, analog für $\cap$. De Morgan: $(A\cup B)^c = A^c\cap B^c$ und $(A\cap B)^c = A^c\cup B^c$.},
type=dinge
}

\newglossaryentry{def_disjunkt_zwei_mengen}{
name={Wann heißen zwei Mengen $A,B$ disjunkt?},
description={Wenn $A\cap B = \emptyset$ gilt.},
type=dinge
}

\newglossaryentry{def_disjunkt_mengenfamilie}{
name={Wann heißt eine Familie $A_1,A_2,\dots$ (allgemein $A_i, i\in I$) disjunkt?},
description={Wenn der Gesamtschnitt aller Mengen leer ist: $\bigcap_i A_i = \emptyset$. (Es reicht, dass es keinen gemeinsamen Punkt in ALLEN Mengen gibt — einzelne Paare können sich trotzdem überschneiden!)},
type=dinge
}

\newglossaryentry{def_paarweise_disjunkt}{
name={Wann heißt eine Familie $A_1,A_2,\dots$ paarweise disjunkt?},
description={Wenn $A_i \cap A_j = \emptyset$ für alle $i\neq j$ gilt (jedes Paar einzeln betrachtet ist disjunkt). Achtung: In der Quelldefinition steht fälschlich $\neq\emptyset$, korrekt ist aber $=\emptyset$ für paarweise Disjunktheit.},
type=dinge
}

\newglossaryentry{def_disjunkte_vereinigung_notation}{
name={Welche Notation verwendet man für die Vereinigung paarweise disjunkter Mengen?},
description={$\dot\bigcup_i A_i$ statt $\bigcup_i A_i$, wenn die $A_i$ paarweise disjunkt sind. Man spricht von einer disjunkten Vereinigung.},
type=dinge
}

\newglossaryentry{bem_paarweise_disjunkt_impliziert_disjunkt}{
name={Folgt aus paarweiser Disjunktheit die (einfache) Disjunktheit einer Mengenfamilie?},
description={Ja: paarweise disjunkt $\Rightarrow$ disjunkt, denn $\bigcap_i A_i \subseteq A_1\cap A_2 = \emptyset$.},
type=dinge
}

\newglossaryentry{bsp_disjunkt_wuerfeln}{
name={Beispiel: Sind $A=\{2,4,6\}$, $B=\{3,6\}$, $C=\{1\}$ beim Würfeln disjunkt?},
description={$A\cap B = \{6\} \neq \emptyset \Rightarrow A,B$ NICHT disjunkt. $A\cap C = \emptyset \Rightarrow A,C$ disjunkt. Zweites Beispiel: $A_1=\{1,3\}, A_2=\{3,5\}, A_3=\{1,5\}$ sind disjunkt ($A_1\cap A_2\cap A_3=\emptyset$), aber NICHT paarweise disjunkt (da z.B. $A_1\cap A_2=\{3\}\neq\emptyset$).},
type=dinge
}

\newglossaryentry{lem_disjunkt_nicht_paarweise_disjunkt}{
name={Gilt auch die Umkehrung: disjunkt $\Rightarrow$ paarweise disjunkt?},
description={Nein. Disjunkt impliziert NICHT paarweise disjunkt — es gibt Gegenbeispiele (siehe $A_k=\{k,k+1,\dots\}$: die Mengen sind disjunkt, aber keine zwei von ihnen sind paarweise disjunkt, da $A_i\cap A_j = A_j\neq\emptyset$ für $i<j$).},
type=dinge
}

\newglossaryentry{bsp_disjunkt_nicht_paarweise_Ak}{
name={Beispiel: Warum sind $A_k=\{k,k+1,k+2,\dots\}$ für $k\in\mathbb{N}$ disjunkt, aber nicht paarweise disjunkt?},
description={Nicht paarweise disjunkt: Für $i<j$ gilt $A_i\cap A_j = A_j \neq \emptyset$. Aber disjunkt: $\bigcap_{i\in\mathbb{N}} A_i=\emptyset$, denn gäbe es ein $\ell$ im Gesamtschnitt, müsste $\ell\in A_{\ell+1}=\{\ell+1,\ell+2,\dots\}$ gelten — unmöglich, da $\ell<\ell+1$.},
type=dinge
}

\newglossaryentry{satz_additivitaet}{
name={Additivität von $P$: Wie berechnet man $P$ einer disjunkten Vereinigung $\dot\bigcup_i A_i$?},
description={$P\left(\dot\bigcup_i A_i\right) = \sum_i P(A_i)$, für paarweise disjunkte $A_1,A_2,\dots$. Gilt auch bei durch Dichten gegebenem $P$. Beweis: Umsortieren der Doppelsumme $\sum_i \sum_{\omega\in A_i}\mu(\omega) = \sum_i P(A_i)$.},
type=dinge
}

\newglossaryentry{satz_komplementaerwahrscheinlichkeit}{
name={Wie lautet die Formel für die Komplementärwahrscheinlichkeit (Gegenwahrscheinlichkeit)?},
description={$P(A^c) = P(\Omega\setminus A) = 1-P(A)$. Beweis: $\Omega = A\ \dot\cup\ A^c \Rightarrow 1=P(\Omega)=P(A)+P(A^c) \Rightarrow P(A^c)=1-P(A)$ (folgt aus Additivität).},
type=dinge
}

\newglossaryentry{satz_monotonie}{
name={Monotonie von $P$: Was folgt aus $A\subseteq B$ für $P(A)$ und $P(B)$?},
description={$A\subseteq B \Rightarrow P(A)\leq P(B)$. Beweis: $B = A\ \dot\cup\ (B\setminus A) \Rightarrow P(B) = P(A) + \underbrace{P(B\setminus A)}_{\geq 0} \geq P(A)$.},
type=dinge
}

\newglossaryentry{satz_differenz_teilmenge}{
name={Wie berechnet man $P(A)$, wenn $A\subseteq B$ bekannt ist (Formel mit $B\setminus A$)?},
description={$A\subseteq B \Rightarrow P(A) = P(B) - P(B\setminus A)$.},
type=dinge
}

\newglossaryentry{satz_vereinigung_zwei_mengen}{
name={Wie lautet die Formel für $P(A\cup B)$ bei nicht notwendig disjunkten $A,B$?},
description={$P(A\cup B) = P(A) + P(B) - P(A\cap B)$. Beweis: $A\cup B = A\ \dot\cup\ [B\setminus(A\cap B)]$, dann Additivität und $P(B\setminus(A\cap B)) = P(B)-P(A\cap B)$ einsetzen.},
type=dinge
}

\newglossaryentry{bem_zusammenfassung_vereinigungsregeln}{
name={Wie berechnet man $P(A\cup B)$ bzw. $P(A\cup B\cup C)$ je nachdem ob die Mengen (paarweise) disjunkt sind?},
description={$A,B$ disjunkt: $P(A\dot\cup B)=P(A)+P(B)$. $A,B$ nicht disjunkt: $P(A\cup B)=P(A)+P(B)-P(A\cap B)$. $A,B,C$ paarweise disjunkt: $P(A\cup B\cup C)=P(A)+P(B)+P(C)$. Allgemein (nicht paarweise disjunkt): siehe Ein- und Ausschlussprinzip für drei Ereignisse.},
type=dinge
}

\newglossaryentry{satz_einschluss_ausschluss_drei}{
name={Wie lautet das Ein- und Ausschlussprinzip für drei Ereignisse $A,B,C$?},
description={$P(A\cup B\cup C) = \underbrace{P(A)+P(B)+P(C)}_{\text{Mengen}} - \underbrace{P(A\cap B)-P(A\cap C)-P(B\cap C)}_{\text{Paare}} + \underbrace{P(A\cap B\cap C)}_{\text{Tripel}}$.},
type=dinge
}



\newglossaryentry{satz_einschluss_ausschluss_n}{
name={Wie lautet das allgemeine Ein- und Ausschlussprinzip für $n$ Ereignisse?},
description={$P\left(\bigcup_{i=1}^n A_i\right) = \sum_{k=1}^n (-1)^{k+1} \sum_{1\leq i_1<\dots<i_k\leq n} P(A_{i_1}\cap\dots\cap A_{i_k})$. Man summiert über alle möglichen Schnitte von $k$ Mengen (mit alternierendem Vorzeichen), für $k=1,\dots,n$. Beweis für allgemeines $n$: vollständige Induktion (Fälle $n=2,3$ als Basis).},
type=dinge
}